\section{Correspondencia con la propuesta inicial}
\label{sec:estructura-propuesta}

La propuesta aprobada, titulada \emph{Phishing. Técnicas y métodos de ataque y cómo detectarlos}, contemplaba una estructura modificable. La tabla~\ref{tab:estructura-propuesta} conserva los trece bloques y su orden, según el apartado 3.1 de la propuesta (páginas 5--6), e indica su ubicación o tratamiento final. El cuerpo se organiza en siete capítulos. El análisis de técnicas de ataque se integra en el marco teórico y las técnicas de defensa se distribuyen entre el estado del arte y la evaluación del prototipo, distinguiendo la evidencia publicada por terceros de los resultados de este trabajo.

\begin{table}[H]
\centering\small
\caption{Correspondencia entre la estructura prevista y la memoria final.}
\label{tab:estructura-propuesta}
\begin{tabularx}{\textwidth}{|>{\RaggedRight\arraybackslash}p{5.2cm}|>{\RaggedRight\arraybackslash}X|}
\hline
\rowcolor{uemcgreen}\color{white}\bfseries Bloque de la propuesta & \color{white}\bfseries Ubicación y tratamiento final \\\hline
1. Portada & Portada con el título aprobado y los datos de autoría y convocatoria. \\\hline
2. Resumen & Resumen inicial; se añade su versión en inglés. \\\hline
3. Introducción & Capítulo 1: contexto y objetivos; alcance concretado en el capítulo 5. \\\hline
4. Marco teórico & Capítulo 2: concepto, evolución, modalidades, fases e ingeniería social. \\\hline
5. Estado del arte & Capítulo 3: estudios previos, herramientas y enfoques de detección y prevención. \\\hline
6. Metodología y planificación & Capítulo 4: enfoque, tecnologías, flujo de trabajo y planificación. \\\hline
7. Análisis de técnicas de ataque & Integrado en el capítulo 2: tipos, ejemplos y estrategias de ingeniería social. \\\hline
8. Análisis y evaluación de técnicas de defensa & Capítulo 3: comparación bibliográfica de métodos, eficacia y límites; capítulo 6: evaluación propia. \\\hline
9. Diseño y desarrollo de medidas defensivas & Capítulo 5: arquitectura, componentes, reglas e interfaz; ejemplos de uso en el capítulo 6. \\\hline
10. Pruebas, resultados y discusión & Capítulo 6: ejemplos, resultados y límites. La comparación con soluciones existentes se trata bibliográficamente en el capítulo 3. \\\hline
11. Conclusiones & Capítulo 7: logros, cumplimiento parcial y trabajo futuro. \\\hline
12. Bibliografía y anexos & Referencias y anexos A--E: repositorio, extractos de código, uso, lista de verificación y evidencia. Capturas y tablas en los capítulos 5 y 6. \\\hline
13. Glosario & Sección final situada antes de las referencias para facilitar la consulta. \\\hline
\end{tabularx}
\end{table}

Esta reorganización evita duplicar la explicación de los ataques y permite seguir el recorrido desde la amenaza hasta el prototipo y sus resultados. La comparación experimental se limita a los modos heurístico, neuronal y combinado del sistema desarrollado. Las soluciones ajenas se comparan mediante bibliografía; no se ha realizado un benchmark directo de productos comerciales. La reducción de alcance funcional, distinta del cambio de organización, se detalla en el apartado~\ref{sec:objetivos-propuesta}.
